% Reach for these first. They are the deck's own vocabulary, so a whole
% talk stays consistent and every one survives greyscale printing.

\term{a primary term}        % accent, bold
\altterm{a secondary term}   % teal, italic
\hl{highlighted}             % amber wash behind the text
\ul{underlined}              % plain underline
\alert{alerted}              % beamer's own, mapped to the accent

% --- Arbitrary colour, still from the palette ----------------------
\tint{accent}{text}   % forest green -- structure and emphasis
\tint{gold}{text}     % amber        -- rules and second-level marks
\tint{muted}{text}    % grey-green   -- deprioritised text
\tint{ink}{text}      % near-black   -- body
\tint{termA}{text}    % forest       -- first defined term
\tint{termB}{text}    % teal         -- second defined term

% --- Raw colour, if you really need it -----------------------------
{\color{lucidAccent} some words}
\textcolor{lucidGold}{some words}

% Palette names you can use anywhere, including in tikz and tables:
%   lucidInk lucidAccent lucidGold lucidMuted lucidBand lucidWash
%   lucidTermA lucidTermB

% --- Colouring part of a table or figure ---------------------------
Ours & \term{78} \\
\rowcolor{lucidBand} header & row \\
\fill[lucidAccent] (0,0) rectangle (1,2);   % inside tikz

% Restraint is the point: the reference decks use three or four colours
% for ~95% of all text. Bold is about 3%. If everything is emphasised,
% nothing is.
